\subsection*{A closed Weil operator and its omitted Fourier complement}
The entrywise bound above can be strengthened to an operator statement.
Here the closed operator is the canonical semilocal Weil operator; it is
not the arithmetic generator whose realization is conditional in Section~18.
Fix $\lambda>1$, put $L=2\log\lambda$, $t_n=2\pi n/L$, and write
\[
 P_\lambda=\sum_{1<m\le\lambda^2}\frac{\Lambda(m)}{\sqrt m},
 \qquad \rho(y)=\frac{e^{y/2}}{2\sinh y},
 \qquad h_L=\sinh^2(L/4).
\]
For later constants use the complete archimedean diagonal
\begin{align}
 A_n={}&\int_0^L\rho(y)
 \{2(1-y/L)\cos(t_ny)-2e^{-y/2}\}\,dy \notag\\
 &+\log(4\pi)+\gamma_{\rm E}+\log\tanh(L/2).
 \label{eq:v114-arch-diagonal}
\end{align}
The cancellation at zero makes this an ordinary convergent integral.

\begin{proposition}[Bounded commutator and the operator core]
\label{prop:v114-weil-core}
In the normalized logarithmic Fourier basis the canonical semilocal
Weil operator $\mathsf W_\lambda$ is
\begin{equation}\label{eq:v114-weil-decomposition}
 \mathsf W_\lambda=D_d+[M_b,H],\qquad
 H_{nm}=\begin{cases}(n-m)^{-1},&n\ne m,\\0,&n=m,\end{cases}
\end{equation}
where $D_d$ is real diagonal, $M_b$ is multiplication by the real odd
sequence
\begin{align}
 b_n={}&\frac{32Lh_Ln}{L^2+16\pi^2n^2}
 +\frac1\pi\left\{\int_0^L\rho(y)\sin(t_ny)\,dy
 +\sum_{1<m\le\lambda^2}\frac{\Lambda(m)}{\sqrt m}
                         \sin(t_n\log m)\right\},\notag\\
 d_n={}&\frac{32Lh_L(L^2-16\pi^2n^2)}{(L^2+16\pi^2n^2)^2}
 -2\sum_{1<m\le\lambda^2}\frac{\Lambda(m)}{\sqrt m}
        (1-\log m/L)\cos(t_n\log m)-A_n.
 \label{eq:v114-b-d}
\end{align}
In particular the diagonal is not filled in by a divided-difference limit.
One has
\begin{equation}\label{eq:v114-b-bound}
 \|b\|_\infty\le B_\lambda^{\rm mat}
 :=\frac{4h_L+P_\lambda+1+\pi/4}{\pi},\qquad
 \|[M_b,H]\|\le2\pi B_\lambda^{\rm mat}.
\end{equation}
At fixed $\lambda$, $d_n=\log|n|+O_\lambda(1)$ as $|n|\to\infty$.
Consequently
\begin{align}
 \mathcal D(\mathsf W_\lambda)
 &=\left\{x:\sum_n\log^2(2+|n|)|x_n|^2<\infty\right\},\notag\\
 \mathcal D[\mathsf W_\lambda]
 &=\left\{x:\sum_n\log(2+|n|)|x_n|^2<\infty\right\}.
 \label{eq:v114-weil-domains}
\end{align}
Finite Fourier sums form an operator core. The operator is lower bounded
and has compact resolvent.
\end{proposition}
\begin{proof}
The exact CCM entries give $\tau_{nm}=(b_n-b_m)/(n-m)$ for $n\ne m$
and $\tau_{nn}=d_n$; see also \eqref{eq:beta-bcoeff-direct}.
For $n>0$, expansion $\rho(y)=\sum_{k\ge0}e^{-(2k+1/2)y}$ and
$t_nL=2\pi n$ give
\[
 S_n:=\int_0^L\rho(y)\sin(t_ny)\,dy
 =\sum_{k\ge0}\frac{(1-e^{-(2k+1/2)L})t_n}
                         {(2k+1/2)^2+t_n^2}.
\]
Interchange is justified by integrability of
$\rho(y)|\sin(t_ny)|$. The summands are nonnegative. The $k=0$
term is at most $1$. For $k\ge1$ discard the factor
$1-e^{-(2k+1/2)L}\le1$ and compare the decreasing majorant
$t_n/\{(2k+1/2)^2+t_n^2\}$ with an integral; this gives at most
$\pi/4$. Thus
$0\le S_n\le1+\pi/4$. Oddness treats negative $n$, and $S_0=0$.
The pole part of $b_n$ has absolute value at most $4h_L/\pi$;
the prime sine sum is bounded by $P_\lambda$. This proves the first
bound in \eqref{eq:v114-b-bound}.

On $\ell^2(\mathbb Z)$, $H$ is the discrete Hilbert transform.
The Fourier series $\sum_{k\ne0}e^{ik\theta}/k=i(\pi-\theta)$
for $0<\theta<2\pi$ proves $H^*=-H$ and $\|H\|=\pi$.
Hence $[M_b,H]$ is bounded self-adjoint and has the asserted norm.
Moreover
\[
 A_n=A_0-\int_0^L w_L(y)(1-\cos(t_ny))\,dy,
 \qquad w_L(y)=2\rho(y)(1-y/L).
\]
Since $w_L(y)-1/y$ is integrable on $(0,L)$, the last integral is
$\log|t_n|+O_L(1)$. For example, split
$\int_0^{|t_n|L}(1-\cos v)\,dv/v$ at $1$ and integrate its
oscillatory tail by parts. The other terms in $d_n$ are bounded at
fixed $\lambda$, proving its logarithmic asymptotic.

The diagonal operator is self-adjoint on the first domain in
\eqref{eq:v114-weil-domains}, has finite sequences as an operator
core, and has compact resolvent. Bounded self-adjoint perturbation
preserves these statements and gives the displayed form domain.
This identifies the actual physical Weil realization: the finite sums
are a form core for its lower-bounded closed form by
\cite[Propositions~3.2--3.4]{ConnesConsaniMoscovici2025}, and both
forms have exactly the same entries there. Uniqueness of the form
closure proves equality, without imposing a new boundary condition.
\end{proof}

Every periodic BV proxy, including the repaired $p_\lambda$, therefore
belongs to $\mathcal D(\mathsf W_\lambda)$ at fixed $\lambda$:
its Fourier coefficients are $O_\lambda((1+|n|)^{-1})$.
Furthermore $P_Np_\lambda\to p_\lambda$ in this operator graph norm,
by the diagonal description and boundedness of the commutator. This
fixed-window statement supplies no endpoint-normalized uniform rate,
and does not remove the separate sampler graph defect.

\begin{proposition}[An explicit positive high-frequency tail]
\label{prop:v114-weil-tail}
Set $a=\min(1,L)$ and
\begin{align}
 C_\lambda^{\rm tail}
 :={}&|A_0|+4+\pi/2+a+2a/L+2P_\lambda
                   +2\sinh(L/2)-L,\notag\\
 \Gamma_{\lambda,N}
 :={}&\log\frac{2\pi(N+1)a}{L}-C_\lambda^{\rm tail}.
 \label{eq:v114-tail-constant}
\end{align}
For every vector $f$ in the form domain supported on Fourier indices
$|n|>N$,
\begin{equation}\label{eq:v114-tail-positive}
 QW_\lambda(f,f)\ge\Gamma_{\lambda,N}\|f\|_2^2.
\end{equation}
In particular this omitted complement is strictly positive whenever
$\Gamma_{\lambda,N}>0$. The constants use no nontrivial-zero data.
\end{proposition}
\begin{proof}
On $0<y<a$, $2\rho(y)\ge e^{-y/2}/y\ge1/y-1/2$, whence
\[
 w_L(y)\ge1/y-(1/2+1/L).
\]
Since $w_L\ge0$ on the full interval and
$F(T):=\int_0^T(1-\cos v)\,dv/v\ge\log T-2$ for $T>0$,
\[
 -A_n\ge\log(|t_n|a)-|A_0|-2-a-2a/L\qquad(n\ne0).
\]
The bound for $F$ is immediate for $T\le1$; for $T\ge1$ the
integral of $\cos v/v$ on $[1,T]$ has absolute value at most $2$
by integration by parts. The off-diagonal archimedean operator has
norm at most $2+\pi/2$, using $b_n^{\rm arch}=S_n/\pi$ in the
preceding commutator proof.

The complete prime operator is a sum of negative weighted truncated
translations and their adjoints, so its lower bound is $-2P_\lambda I$.
In centered logarithmic coordinates $x\in[-L/2,L/2]$, the pole form is
\[
 2|\ip f{\cosh(x/2)}|^2-2|\ip f{\sinh(x/2)}|^2
 \ge-\{2\sinh(L/2)-L\}\|f\|_2^2.
\]
Combining these three estimates gives
\eqref{eq:v114-tail-positive} first for finite Fourier sums and then
on the tail form domain by closure.
\end{proof}

This is a quantitative omitted-tail theorem, not full Weil positivity.
Its absolute arithmetic bound still gives a sufficient cutoff
exponential in $O(\lambda)$; it does not reach the natural prolate
scale. With $Q_N=I-P_N$ and $\Gamma_{\lambda,N}>0$, write
\[
 \mathsf W_\lambda=
 \begin{pmatrix}F&B^*\\B&T\end{pmatrix},\qquad
 F=P_N\mathsf W_\lambda P_N,\quad
 B=Q_N\mathsf W_\lambda P_N,\quad T\succeq\Gamma_{\lambda,N}I.
\]
The cross operator $B$ is bounded by \eqref{eq:v114-weil-decomposition}.
The remaining exact sign problem is the finite matrix
$K=F-B^*T^{-1}B$. There is a useful verified-solve alternative to
bounding $B$ alone. For any $Z:E_N\to\mathcal D(T)$ put $R=B-TZ$.
Then
\begin{equation}\label{eq:v114-schur-residual}
 K=F-B^*Z-Z^*B+Z^*TZ-R^*T^{-1}R
 \succeq F-B^*Z-Z^*B+Z^*TZ
             -\frac{\|R\|^2}{\Gamma_{\lambda,N}}I.
\end{equation}
This follows by expansion. If the final finite lower bound is
$-\varepsilon I$, $\varepsilon\ge0$, completing the form square
gives $\mathsf W_\lambda\succeq-\varepsilon I$.
Proving such a signed bound along growing windows is still open;
positivity of $F$ alone is insufficient.
